\section{第\ref{sec:code}章　モールス符号}

この章の上限の見積もりは情報理論と計算によるものである。測定されているのは、通信路のどこでノイズが生じるか、脳がどれだけ速く動作するか、スパイクがどれだけのコストをもつかである。皮質が実際にどの符号を使っているかという問いは未解決のままである。

{\footnotesize
\setlength{\tabcolsep}{3pt}\renewcommand{\arraystretch}{1.15}
\begin{longtable}{>{\raggedright}p{0.5cm}>{\raggedright}p{4.0cm}>{\raggedright}p{5.6cm}>{\raggedright}p{2.3cm}>{\raggedright\arraybackslash}p{1.7cm}}
\toprule
No. & 主張 & 実験と測定されたもの & 出典 & 状態 \\
\midrule\endfirsthead
\toprule
No. & 主張 & 実験と測定されたもの & 出典 & 状態 \\
\midrule\endhead
1 & 皮質の\nt{GLUT}ニューロンの発火率は毎秒1スパイク未満から数十スパイクで、ニューロンはほとんどの時間を下限の近くで動作する & 覚醒ラットの聴覚野で、細胞に密着させたピペットによる記録（ニューロンの選択はその活動によらない）：各瞬間に毎秒$20$スパイクを超える発火率を示すニューロンは$5\,\%$未満。一部のニューロンの背景発火率は毎秒$0.01$未満。発火率の分布は対数正規 & \cite{Hromadka2008} & 測定済み \\
2 & スパイク通信路の容量は$R_{\max}\approx r\log_2\frac{e}{r\Delta t}$~\eqref{eq:spikeentropy}で、そのほとんどすべてがスパイクの時刻にある（命題~\ref{prop:capacity}） & 長さ$\Delta t$のセルからなる2進文字列のエントロピー。章の数値：$r=10$スパイク/s、$\Delta t=1\ms$でスパイク1個あたり$8$ビット、$\Delta t=10\ms$で約$5$ビット、スパイクカウンタでは毎秒$2$ビット未満 & \cite{MacKay1952} & 理論 \\
3 & 感覚ニューロンはスパイク1個あたり1ビットから数ビットを伝え、最もよい場合で上限の半分に達する & 刺激がわかっている感覚ニューロンのスパイクから刺激を復元。測定のまとめはこの書籍にある & \cite{Rieke1997} & 測定済み \\
4 & スパイク発生器は精確である。精確な時刻を決めるのは入力の速い変化である & ラット皮質のスライスのニューロン：速く揺らぐ電流を繰り返し与えるとスパイクは$1\ms$未満の精度で再現され、一定の電流では時刻がばらける & \cite{Mainen1995} & 測定済み \\
5 & シナプスは当てにならない。放出のランダム性のため、スパイクの半分以上は受け手に届かない & 海馬スライスで錐体ニューロンの入力を最小刺激：スパイクの半分以上が応答を起こさない。刺激閾値の揺らぎ、軸索分岐での伝導の失敗、実験温度の低さは除外された。残るのは確率的な放出である & \cite{AllenStevens1994} & 測定済み \\
6 & 生きた皮質のスパイク列はランダムに見える。間隔はポアソン列のようにばらつき、スパイク数の分散は平均に近い。「ノイズ」の一部は動物の運動についてのメッセージである & 覚醒サルの視覚野V1とMTでの記録：間隔の変動係数は約$1$、スパイク数の分散と平均の比はランダム過程と同様。独立な小さな入力を足し合わせるモデルでは、ばらつきははるかに小さくなる。マウス視覚野では、ばらつきのかなりの部分が動物自身の運動のビデオ記録から予測できる & \cite{Softky1993,Stringer2019} & 測定済み \\
7 & レート符号は遅い。誤差は$1/\sqrt{rT}$~\eqref{eq:rateerror}だが、脳は画像を$\sim150\ms$で認識する & 式はポアソン統計で、章の導出。実験：$20\ms$だけ見せた見知らぬ写真に動物がいるかどうかを被験者が判断した。「はい」と「いいえ」の脳電位は約$150\ms$後に分かれる。この時間を$10$--$20\ms$ずつ約10段に分けるのは章の推定であり、測定ではない & \cite{Thorpe1996} & 測定済み \\
8 & 群での平均化は相関によって制限される。誤差は$1/\sqrt c$分の1より小さくならない~\eqref{eq:correlated}（命題~\ref{prop:pooling}） & 同時記録したサルMT野のニューロンのペア：スパイク数の相関係数は平均$0.12$で、理論解析によればこの程度の相関でも群による改善は制限される。理論：限界を決めるのは相関のうち信号に似た部分だけである。反対の立場のモデル：$50$--$100$個のニューロンの群の発火率は1スパイク間隔、$10$--$50\ms$で読まれ、個々のスパイクの時刻は皮質では使われない & \cite{Zohary1994,MorenoBote2014,ShadlenNewsome1998} & 測定済み \\
9 & 数は最初のスパイクの時刻~\eqref{eq:latency}、群のスパイクの順序、局所リズムの位相に書くことができる & サンショウウオの網膜：短時間提示した画像の空間構造は、出力ニューロンの最初のスパイクの相対遅延に書かれており、コントラストにほとんどよらない。ラット海馬の場所細胞：「自分の」場所を通過するにつれて、スパイクはシータリズム（$7$--$12$\,Hz）のサイクルのますます早い時点へずれ、ずれは$100^\circ$から$355^\circ$に及ぶ。皮質がスパイクの時刻をどれだけ使うかは未解決の問題（8行目） & \cite{Gollisch2008,OKeefe1993} & 測定済み \\
10 & どの符号が読まれるかは受け手の時定数が決める~\eqref{eq:reader}（命題~\ref{prop:reader}）。活動中の皮質ではニューロンは同時検出器に近い & 式~\eqref{eq:reader}は章の導出。in vivo記録の総説：活動中の皮質では膜のコンダクタンスが静止時の数倍になり、時定数はそれに応じて短い。皮質がまさにこのように符号を切り替えることは直接には示されていない & \cite{Destexhe2003} & 仮説 \\
11 & 枯渇するシナプスは発火率の変化、本質的には$\Delta r/r$を伝える~\eqref{eq:depression}（図~\ref{fig:depression}） & 皮質錐体ニューロンのペア記録：枯渇の速さは放出確率で決まる。枯渇が遅いと応答は発火率を反映し、速いと同時性を反映する。ラット皮質で測定された枯渇に基づくモデル：速い入力と遅い入力で発火率の相対変化が等しければ応答も等しく、つまりシナプスは$\Delta r/r$を伝える。数値$1.7\to2.2$、$6.7$、$90\ms$は章の計算 & \cite{Tsodyks1997,Abbott1997depression} & 測定済み \\
12 & バーストは「ツー」である。失われる確率$(1-p)^k$~\eqref{eq:burst}は小さく、促通がそれをさらに下げる & 総説：当てにならないシナプスでは単発のスパイクはしばしば失われるが、バーストは促通のおかげで確実に伝わる。バーストはシナプスの長期的な変化を起こす。脳がバーストを独立した記号として読むというのは仮説 & \cite{Lisman1997,AllenStevens1994} & 仮説 \\
13 & 速い\nt{GABA}ニューロンはリズムと「句読点」を与える。もう一つのクラスは抑制するものを抑制し、ゲートを開く（命題~\ref{prop:punctuation}） & 皮質の\nt{GABA}ニューロンのクラスの性質の総説。光遺伝学：速い\nt{GABA}ニューロンのリズミックな興奮はガンマリズムを起こし、刺激への応答はサイクルの位相に依存する。マウス視覚野では、PVニューロンは互いを抑制し、SSTニューロンは他のすべてのクラスを、VIPニューロンはおもにSSTを抑制する。覚醒マウスでVIPニューロンを興奮させると\nt{GLUT}ニューロンの抑制が解除される。VIPニューロンをオンにするのは強化の信号である。一般規則としての役割分担は仮説 & \cite{Tremblay2016,Cardin2009,Pfeffer2013,Pi2013} & 仮説 \\
14 & エネルギーは平均発火率を毎秒1スパイク程度に制限する~\eqref{eq:energy} & 解剖学的・生理学的データによるげっ歯類灰白質のエネルギー収支：スパイクとシナプス電流は信号伝達のエネルギーの$47\,\%$と$34\,\%$。同時に活動できるニューロンは$\sim15\,\%$以下。ヒトの皮質では、同時に強く活動できるニューロンは$\sim1\,\%$以下。スパイク1個あたり$\sim10^9$個のATPと信号伝達の電力$\sim10$\,Wは、これらの計算に基づく章の推定 & \cite{Attwell2001,Lennie2003} & 理論 \\
15 & 符号はスパースで、1ジュールあたりのビット数が最大になるよう調整されている。皮質は精度をエネルギーと引き換えにする（命題~\ref{prop:economy}） & 経済的な符号の仮説。根拠：食餌制限したマウスの視覚野では、AMPA受容体のコンダクタンスとシナプスのATP消費が下がり（$-29\,\%$）、スパイクの発火率は保たれ、方位選択性は$32\,\%$広がる & \cite{Barlow1961,Padamsey2022} & 仮説 \\
\bottomrule
\end{longtable}}
